\documentclass[11pt]{article}
\usepackage[margin=1in]{geometry}
\usepackage{amsmath,amssymb,amsthm}
\usepackage{array}
\usepackage{parskip}

\newtheorem{theorem}{Theorem}
\newtheorem{lemma}[theorem]{Lemma}
\newtheorem{corollary}[theorem]{Corollary}
\newtheorem{proposition}[theorem]{Proposition}
\theoremstyle{definition}
\newtheorem{definition}[theorem]{Definition}
\theoremstyle{remark}
\newtheorem{remark}[theorem]{Remark}

\newcommand{\swapgraph}{\mathrm{swap\_graph}}
\DeclareMathOperator{\dist}{dist}

\title{Disproof of conjecture \texttt{00000008234}}
\author{earthking11}
\date{2026-09-13}

\begin{document}
\maketitle

\section*{Verdict: FALSE}

We disprove conjecture \texttt{00000008234} as stated. The conjecture is a
conjunction of several clauses, the first of which asserts that for an arbitrary
number of agents dividing an arbitrary number of indivisible goods, the
EF1 exchange graph is connected with diameter at most $m(n-1)$ under arbitrary
monotone valuations. We exhibit an elementary counterexample with $n = 2$ agents,
$m = 2$ goods and unit valuations: the EF1 exchange graph has exactly two
vertices and \emph{no} edge, so it is disconnected and its diameter is infinite.
The claimed bound $m(n-1) = 2$ therefore fails, and the first conjunct of the
conjecture is false; hence the whole conjunction is false.

\section{The conjecture}

Quoted verbatim from \texttt{conjectures/00000008234.md}:

\begin{quote}
\textbf{Definition:} The exchange graph $\swapgraph$ of envy-free-up-to-one-good
allocations when $n$ agents divide $m$ indivisible goods: vertices are EF1
allocations and edges are single-good transfers.
\textbf{Conjecture:} Under arbitrary monotone valuations $\swapgraph$ is
connected with diameter at most $m(n-1)$; an allocation that is simultaneously
EF1 and Pareto optimal always exists and is reachable by a polynomial-time
item-moving algorithm; and EFX allocations always exist for three agents, while
a minimal counterexample, if one exists, must have exactly 4 agents and require a
non-bimodal valuation profile.
\end{quote}

We refute the first conjunct. Throughout, $\swapgraph$ has as vertices the EF1
allocations and as edges the unordered pairs of EF1 allocations that differ by a
single-good transfer.

\section{Allocations, unit valuations, and EF1}

\begin{definition}[allocation]
An allocation of $m$ indivisible goods to $n$ agents is a map
$a : \{1,\dots,m\} \to \{1,\dots,n\}$ sending each good to its owner. The bundle
of agent $i$ is $A_i = a^{-1}(i)$.
\end{definition}

\begin{definition}[unit additive valuation]
The \emph{unit} valuation assigns every agent the same value $v_i(g) = 1$ to
every good $g$, extended additively to bundles: $v_i(A) = |A \cap A_i|$ for the
bundle $A_i$ of agent $i$. Unit valuations are additive, hence monotone: if
$A \subseteq B$ then $v_i(A) \le v_i(B)$. Thus they are admissible under the
conjecture's hypothesis of arbitrary monotone valuations.
\end{definition}

\begin{definition}[EF1]
An allocation $a$ is \emph{envy-free up to one good} (EF1) if for every ordered
pair of agents $(i,j)$ there is a good $g \in A_j$ such that
\[
v_i(A_j \setminus \{g\}) \;\le\; v_i(A_i).
\]
Equivalently, whenever $i$ envies $j$ (that is, $v_i(A_j) > v_i(A_i)$), the
removal of a single suitable good from $A_j$ eliminates the envy.
\end{definition}

\begin{lemma}[unit-valuation characterisation]\label{lem:gap}
Under unit valuations, an allocation $a$ is EF1 if and only if
\[
\max_i |A_i| - \min_i |A_i| \;\le\; 1 .
\]
\end{lemma}

\begin{proof}
Write $s_i = v_i(A_i) = |A_i|$.
($\Leftarrow$) Let $i,j$ be agents. If $s_i \ge s_j$ then $i$ does not envy $j$,
and the EF1 condition is satisfied. Otherwise $s_j \ge s_i + 1$. If
$s_j = s_i + 1$ then $A_j \neq \emptyset$ and removing any $g \in A_j$ gives
$v_i(A_j \setminus \{g\}) = s_j - 1 = s_i \le s_i$, so the condition holds. If
$s_j \le s_i + 1$ for all $i,j$, one of these two cases applies to every ordered
pair.
($\Rightarrow$) Suppose $s_j \ge s_i + 2$ for some $i,j$. For every good
$g \in A_j$ we have $v_i(A_j \setminus \{g\}) = s_j - 1 \ge s_i + 1 > s_i$, so
no single removal eliminates the envy of $i$ towards $j$. Hence $a$ is not EF1.
Contrapositively, EF1 forces $s_j \le s_i + 1$ for all ordered pairs, i.e.
$\max_i s_i - \min_i s_i \le 1$.
\end{proof}

\section{The counterexample: $n = 2$ agents, $m = 2$ goods}

Take goods $\{g_1,g_2\}$ and agents $\{0,1\}$ with unit valuations. There are
$2^2 = 4$ allocations. By Lemma~\ref{lem:gap}, EF1 is equivalent to the bundle
sizes differing by at most $1$. The complete table is:

\begin{center}
\begin{tabular}{c|c|c|c}
allocation $a$ & bundles $(A_0, A_1)$ & sizes $(s_0, s_1)$ & EF1? \\
\hline
both goods to agent $0$ & $(\{g_1,g_2\},\ \emptyset)$ & $(2,0)$ & no \\
$g_1 \mapsto 0,\ g_2 \mapsto 1$ & $(\{g_1\},\ \{g_2\})$ & $(1,1)$ & \textbf{yes} \\
$g_1 \mapsto 1,\ g_2 \mapsto 0$ & $(\{g_2\},\ \{g_1\})$ & $(1,1)$ & \textbf{yes} \\
both goods to agent $1$ & $(\emptyset,\ \{g_1,g_2\})$ & $(0,2)$ & no \\
\end{tabular}
\end{center}

\begin{theorem}\label{thm:two-ef1}
For $n = m = 2$ with unit valuations there are exactly two EF1 allocations,
namely the two diagonal allocations
\[
a_{01} : g_1 \mapsto 0,\ g_2 \mapsto 1,
\qquad
a_{10} : g_1 \mapsto 1,\ g_2 \mapsto 0,
\]
and no others.
\end{theorem}

\begin{proof}
By the table, the allocations with size vector $(2,0)$ and $(0,2)$ have a gap of
$2$ and are not EF1 by Lemma~\ref{lem:gap}. The two allocations with size vector
$(1,1)$ have gap $0 \le 1$ and are EF1. The four allocations are all of them,
since each of the two goods has two possible owners. Hence exactly $a_{01}$ and
$a_{10}$ are EF1.
\end{proof}

\begin{theorem}\label{thm:no-edge}
In the above instance, $\swapgraph$ has no edge. Consequently $\swapgraph$ is
disconnected and has infinite diameter; in particular the conjectured bound
$m(n-1) = 2$ fails.
\end{theorem}

\begin{proof}
The vertices of $\swapgraph$ are $a_{01}$ and $a_{10}$ by
Theorem~\ref{thm:two-ef1}. They differ in both coordinates, so their Hamming
distance is $2$; they are not related by a single-good transfer.

It remains to check that a single-good transfer out of either vertex does not
land on a vertex. Starting from $a_{01}$, transferring $g_1$ to agent $1$ yields
the allocation with sizes $(0,2)$, and transferring $g_2$ to agent $0$ yields the
allocation with sizes $(2,0)$; both have size gap $2$ and are not EF1 by
Lemma~\ref{lem:gap}. The same two targets arise from $a_{10}$. Hence no
single-good transfer has both endpoints EF1, i.e.\ $\swapgraph$ has no edge.

A graph on two vertices with no edge has two connected components, so it is not
connected, and its diameter is infinite (there is no path, so no finite bound on
path lengths). Since the conjecture asserts connectivity with diameter at most
$m(n-1) = 2$, its first conjunct is false.
\end{proof}

\begin{remark}
The argument is even stronger than disconnection: \emph{every} single-good
transfer out of an EF1 allocation leaves the EF1 set. Thus $\swapgraph$ is a
totally disconnected graph on two vertices, and the same happens in every
instance of the next section.
\end{remark}

\section{An infinite family: $n$ divides $m$}

The phenomenon is not an artifact of the smallest instance.

\begin{theorem}\label{thm:family}
Let $n \ge 2$ and $m \ge n$ with $n \mid m$, and use unit valuations. Then every
EF1 allocation is balanced, i.e.\ $|A_i| = m/n$ for every agent $i$. Consequently
every single-good transfer produces an allocation that is not EF1, so
$\swapgraph$ has no edges at all. Since there are at least two balanced
allocations, $\swapgraph$ is disconnected and its diameter is infinite.
\end{theorem}

\begin{proof}
Let $a$ be EF1 and put $s_i = |A_i|$. By Lemma~\ref{lem:gap},
$\max_i s_i - \min_i s_i \le 1$, so there is an integer $k$ with
$s_i \in \{k, k+1\}$ for all $i$. Let $r$ be the number of agents with
$s_i = k+1$; then
\[
m = \sum_{i=1}^{n} s_i = nk + r, \qquad 0 \le r \le n .
\]
Since $n \mid m$, we have $r \equiv 0 \pmod n$, so $r = 0$ or $r = n$, and in
both cases all $s_i$ are equal, with common value $m/n$. (Equivalently: an
average that is an integer and all sizes within $1$ of one another forces all
sizes to equal the average.)

Now perform one single-good transfer: some agent $i$ loses a good and some agent
$j$ gains one, so the new sizes are $m/n - 1$ for $i$ and $m/n + 1$ for $j$,
while every other agent still has $m/n$. The gap is
$(m/n + 1) - (m/n - 1) = 2 > 1$, so by Lemma~\ref{lem:gap} the new allocation is
not EF1. As every EF1 allocation has every agent at exactly $m/n$ goods, no
single-good transfer has both endpoints EF1, and $\swapgraph$ has no edges.

Finally, the number of balanced allocations is the multinomial
$m!/(\,(m/n)!\,^{n})$, which is at least $2$ when $n \ge 2$ and $m \ge n$
(for $m = n$ it is $n! \ge 2$, and for $m/n \ge 2$ there are again at least two
ways to choose one agent's share). With at least two vertices and no edge,
$\swapgraph$ is disconnected.
\end{proof}

\begin{corollary}[$n = m = 3$]
Three agents and three goods with unit valuations form a counterexample with
three (not two) agents. The EF1 allocations are exactly the six permutations
that give each agent one good; all pairwise Hamming distances are at least $2$,
so $\swapgraph$ has no edges and is disconnected.
\end{corollary}

\begin{proof}
By Theorem~\ref{thm:family} with $n = m = 3$, every EF1 allocation assigns
exactly $3/3 = 1$ good to each agent. There are $3! = 6$ such allocations. Two
distinct permutations differ in at least two positions, so their Hamming
distance is at least $2$; a single-good transfer moves exactly one good, landing
on a size vector with a gap of $2$, which is not EF1. Hence no edge, and the six
vertices split into six components.
\end{proof}

The following table records the disconnection for several members of the family
(verified by exhaustive enumeration in \texttt{reproduce.py}):

\begin{center}
\begin{tabular}{c|c|c|c}
$n$ & $m$ & \# EF1 allocations & \# edges \\
\hline
$2$ & $2$ & $2$ & $0$ \\
$2$ & $4$ & $6$ & $0$ \\
$2$ & $6$ & $20$ & $0$ \\
$3$ & $3$ & $6$ & $0$ \\
$4$ & $4$ & $24$ & $0$ \\
\end{tabular}
\end{center}

\section{Conclusion}

With $n = 2$ agents, $m = 2$ goods and unit (hence monotone, additive)
valuations, the EF1 exchange graph $\swapgraph$ has exactly two vertices and no
edges. It is disconnected, its diameter is infinite, and the conjectured bound
$m(n-1) = 2$ fails. The same holds for every $n \mid m$ with $n \ge 2$. This
refutes the first conjunct of conjecture \texttt{00000008234}, and therefore the
conjecture as a conjunction.

\section{Reproducibility and formalisation}

The exhaustive enumeration, the literal EF1 predicate, the swap graph and its
connected components are computed by \texttt{reproduce.py} (standard-library
Python~3 only), which prints the tables above, asserts zero edges and
disconnection, and ends with \texttt{PASS}.

The $2 \times 2$ refutation is formalised in Lean~4 in the accompanying
\texttt{lean4/} project, core Lean only (no Mathlib, no \texttt{sorry}). The main
formal statements are

\begin{quote}\ttfamily
theorem exactly\_two\_ef1 : ef1 alloc01 \(\wedge\) ef1 alloc10 \(\wedge\) \dots\\
theorem no\_adj\_ef1 : ef1 a \(\to\) ef1 b \(\to\) \(\neg\) Adj a b\\
theorem not\_ef1\_connected : \(\neg\) EF1Connected\\
theorem no\_path\_between\_diagonals : \(\neg\) (\(\exists\) k, ReachN alloc01 k alloc10)\\
theorem claimed\_diameter\_bound\_fails : \(\neg\) SwapGraphDiameterAtMost 2
\end{quote}

\noindent Here \texttt{ef1} is the EF1 predicate transcribed literally (for
every ordered pair of agents, some good can be removed from the envied bundle to
remove the envy), \texttt{Adj} is Hamming distance $1$ (a single-good transfer),
\texttt{ReachN} is reachability through EF1 vertices, and
\texttt{SwapGraphDiameterAtMost} is the conjectured connectivity plus diameter
bound. \texttt{Check.lean} prints the axioms of each theorem, confirming that
only \texttt{propext} and \texttt{Quot.sound} are used and that \texttt{sorryAx}
does not occur.

\end{document}
